\begin{textsample}{POS Dim 3 – 2009\_c4 – Score 1.00 – 2009\_c4/enade\_2009\_arquivol\_ba1f6e9589.txt}  \label{ex:f3_pos_122}
Capítulo 4 Percepção sobre a Prova As percepções dos estudantes sobre a prova aplicada no ENADE/2009 na área de ARQUIVOLOGIA foram mensuradas por meio de nove itens que avaliaram desde o grau de dificuldade da prova até o tempo gasto para concluí -la.
As questões foram analisadas separando -se concluintes e ingressantes, e as percepções sobre a prova foram relacionadas com o desempenho dos estudantes e com a Região Geográfica de localização do curso6.
O desempenho dos estudantes dentro de cada Área/Subárea foi classificado em quartos ( 1/4 ) para os concluintes e para os ingressantes em separado ( ver Figura 1 para o esquema de representação gráfica destas grandezas ).
As notas finais foram colocadas em ordem ascendente e no primeiro quarto foram colocadas as 25 \% notas mais baixas, limitadas superiormente pelo percentil 25 que vem a ser o primeiro quartil.
No segundo quarto, as 25 \% seguintes, limitadas superiormente pelo percentil 50 que vem a ser a mediana.
No terceiro quarto, as 25 \% notas após as encontradas no segundo quarto, ou seja, as primeiras 25 \% acima da mediana, limitadas superiormente pelo percentil 75 que vem a ser o terceiro quartil.
No último quarto foram alocadas as 25 \% maiores notas.
Potencialmente, poderia existir uma relação entre a nota do estudante e a sua percepção sobre a prova.
Calculando as médias das respostas segundo os quartos de notas, seria possível reconhecer a existência ( ou não ) de algum gradiente nas respostas, por exemplo, uma melhor nota correspondendo a uma percepção menos negativa.
Este gradiente seria traduzido, por exemplo, no gráfico 4.1a, pela diminuição das áreas \textbf{laranja} e vermelha ( Difícil e Muito difícil ) com o aumento do desempenho ( indo do primeiro ao último quarto ). 6 Curso, neste relatório, refere -se à unidade de análise para o Conceito ENADE e é caracterizado pela combinação de IES, Área/Subárea ( quando pertinente ) e município de habilitação.
A medida tradicional de “ grau de discriminação ” comparou os resultados dos estudantes do primeiro quarto ( 25 \% piores estudantes ) aos do último quarto ( 25 \% melhores ).
Tal comparação, no gráfico disponibilizado, foi idêntica à comparação da primeira barra ( primeiro quarto ) com a última barra ( último quarto ) de cada grupo ingressante ou concluinte nos gráficos “ segundo quartos de nota final ” ( gráficos com números ímpares neste capítulo ).
Por outro lado, a disponibilização das informações segundo os quartos de nota permitiu verificar o comportamento de dois outros grupos intermediários que na análise somente dos extremos tem que ser suposto monotônico.
A seguir, serão apresentados os principais resultados relativos aos nove itens avaliados. 4.1 Grau de dificuldade da prova 4.1.1 Formação Geral Os Gráficos 4.1a, 4.1b e 4.1c apresentam a distribuição das respostas à questão “ Qual o grau de dificuldade desta prova na parte de Formação Geral? ”, segundo quartos de nota final e grupo de estudantes concluintes e ingressantes que participaram do ENADE/2009 na área de ARQUIVOLOGIA.
No total de concluintes, 0,9 \% declararam que a prova na parte da Formação Geral foi Muito difícil, contra um valor de 1,9 \% para os ingressantes.
Tanto entre os concluintes, como entre os ingressantes, a proporção de indivíduos que classificaram esta parte como Difícil/Muito difícil diminui com o desempenho do estudante, como era esperado.
Por outro lado, a proporção que declarou a parte Fácil/Muito fácil foi maior entre os concluintes ( 23,3 \% ) do que entre os ingressantes ( 21,5 \% ).
Não se observou uma variação monotônica significativa desta proporção com o aumento no desempenho, nem entre os concluintes, nem entre os ingressantes.
Entre os concluintes, 63,0 \% classificaram esta parte com um nível Médio de dificuldade.
A proporção entre os ingressantes foi ligeiramente maior, 65,1 \%.
Entre os concluintes, esta proporção aumentou com o desempenho, indo de 53,3 \% no primeiro quarto de notas a 68,5 \% no último quarto.
Os Gráficos 4.2a e 4.2b apresentam a distribuição das respostas à questão “ Qual o grau de dificuldade desta prova na parte de Formação Geral? ”, segundo a Região onde o curso está situado para os estudantes concluintes e ingressantes que participaram do ENADE/2009 na área de ARQUIVOLOGIA.
Cumpre lembrar que a Região Norte não apresentou nenhum estudante concluinte.
A Região Sudeste foi a única com declaração de Muito difícil.
Entre os concluintes, a maior proporção de estudantes que considerou a parte de Formação Geral da prova como de dificuldade média entre todas as Regiões brasileiras, foi da Região Sul com 73,6 \%, já entre os ingressantes, foram os estudantes da Região Norte com 76,9 \%.
As Regiões com as menores proporções desta resposta foram respectivamente, para os concluintes e ingressantes, a Centro-Oeste ( 50,0 \% ) e a Sul ( 60,7 \% ). 4.1.2 Componente Específico Os Gráficos 4.3a, 4.3b e 4.3c apresentam a distribuição das respostas à questão “ Qual o grau de dificuldade desta prova na parte do Componente Específico? ”, segundo quartos de nota final e grupo de estudantes concluintes e ingressantes que participaram do ENADE/2009 na área de ARQUIVOLOGIA.
Em geral, para as desagregações consideradas, os estudantes classificaram o Componente Específico como mais difícil que a Formação Geral.
No grupo de concluintes, 18,6 \% declarou que a prova na parte do Componente Específico foi Difícil/Muito difícil, sendo um valor muito maior para os ingressantes, 33,4 \%.
Como esperado, entre os concluintes, a proporção de indivíduos que classificaram esta parte como Muito difícil, diminuiu com o desempenho do estudante.
Entre os ingressantes, esta proporção decresceu com o desempenho até o terceiro quarto.
Entre os concluintes nenhum considerou esta parte Muito fácil.
Entre os ingressantes, somente 2,4 \% dos estudantes no primeiro quarto assim o fizeram.
A proporção que declara a parte Fácil/Muito fácil foi maior entre os concluintes ( 15,0 \% ) do que entre os ingressantes ( 9,0 \% ).
Entre os concluintes, contrariamente ao esperado, a proporção dos que consideraram esta parte como Fácil decresceu com o desempenho.
O comportamento oposto se verificou entre os ingressantes.
Entre os concluintes, 66,4 \% classificaram esta parte com um nível Médio de dificuldade.
A proporção entre os ingressantes foi ligeiramente menor, 57,7 \%.
Entre os concluintes, esta proporção basicamente aumentou com o desempenho, indo de 60,0 \% no primeiro quarto de notas, caindo para 58,5 \% no segundo quarto, chegando a 70,7 \% no último quarto.
Entre os ingressantes, esta proporção foi crescente nos primeiros dois quartos, com uma queda no terceiro.
Os Gráficos 4.4a e 4.4b apresentam a distribuição das respostas à questão “ Qual o grau de dificuldade desta prova na parte do Componente Específico? ”, segundo a Região onde o curso está situado para os estudantes concluintes e ingressantes que participaram do ENADE/2009 na área de ARQUIVOLOGIA.
Cumpre lembrar aqui também que a Região Norte não apresentou nenhum estudante concluinte.
As respostas dadas pelos estudantes concluintes de cada Região mostraram que nenhum concluinte declarou que esta parte é Muito fácil, e que a Região Sudeste foi a única com declaração de Muito difícil.
Entre os concluintes, a maior proporção de estudantes que considerou a parte do Componente Específico da prova como de dificuldade média entre todas as Regiões brasileiras, foi na Região Sul com 81,1 \% ; já entre os ingressantes, foi os estudantes da Região Nordeste com 67,3 \%.
As Regiões com as menores proporções desta resposta foram respectivamente para os concluintes e ingressantes, a Nordeste ( 50,0 \% ) e a Norte ( 20,0 \% ). 4.2 Extensão da prova em relação ao tempo total Os Gráficos 4.5a, 4.5b e 4.5c apresentam a distribuição das respostas à questão sobre a extensão da prova em relação ao tempo total, segundo quartos de nota final e grupo de estudantes concluintes e ingressantes que participaram do ENADE/2009 na área de ARQUIVOLOGIA.
Tanto para concluintes quanto para ingressantes, a maioria considerou a extensão da prova como adequada em relação ao tempo destinado à sua resolução, respectivamente 62,7 \% e 64,1 \%.
Entre os concluintes, esta proporção cresceu com o desempenho nos três primeiros grupos, caindo no último.
Entre os ingressantes, esta proporção foi crescente a partir do segundo grupo, alcançando 69,9 \% para o último quarto.
A proporção dos que consideraram a extensão da prova muito curta foi ligeiramente maior entre os concluintes ( 1,3 \% ) do que entre os ingressantes ( 1,0 \% ).
A proporção dos que consideraram a extensão da prova curta/muito curta quando comparados ao tempo total foi maior entre os concluintes ( 8,0 \% ) do que entre os ingressantes ( 3,7 \% ).
Entre os concluintes, somente estudantes com desempenho no grupo mais baixo declararam que a extensão da prova teria sido muito curta com relação ao tempo total.
Entre os ingressantes, esta declaração só não aconteceu no terceiro quarto.
A maior proporção de concluintes declarando que a prova foi curta/muito curta foi entre os estudantes de pior desempenho, 23,3 \%.
A proporção dos que consideraram a extensão da prova muito longa foi quase idêntica entre os concluintes ( 9,8 \% ) e ingressantes ( 9,7 \% ).
A proporção dos que consideraram a extensão da prova longa foi menor entre os concluintes ( 19,6 \% ) do que entre os ingressantes ( 22,6 \% ).
Os Gráficos 4.6a e 4.6b apresentam a distribuição das respostas à questão sobre a extensão da prova em relação ao tempo total, segundo a Região onde o curso está situado para os estudantes concluintes e ingressantes que participaram do ENADE/2009 na área de ARQUIVOLOGIA.
Em todas as Regiões brasileiras, a maioria dos estudantes, tanto concluintes, quanto ingressantes, percebeu a prova como adequada. 4.3 Compreensão dos Enunciados das Questões 4.3.1 Formação Geral O Gráfico 4.7a, 4.7b e 4.7c apresentam a distribuição das respostas à questão “ Os enunciados das questões da prova na parte de Formação Geral, estavam claros e objetivos? ”, segundo quartos de nota final e grupo de estudantes concluintes e ingressantes que participaram do ENADE/2009 na área de ARQUIVOLOGIA.
Com relação aos enunciados das questões da parte de Formação Geral, as opiniões foram positivas : a grande maioria dos concluintes ( 75,8 \% ) e ingressantes ( 79,0 \% ) consideram claros e objetivos os enunciados de todas ou da maioria das questões.
Tanto entre os concluintes quanto entre os ingressantes, esta proporção cresceu com o desempenho.
Apenas dois concluintes e dois ingressantes declararam que não, nenhum dos enunciados das questões da parte de Formação Geral era claro e objetivo, correspondendo a 6,7 \% dos concluintes no primeiro quarto, 1,2 \% dos ingressantes do primeiro quarto e 0,8 \% dos ingressantes do segundo quarto.
Em linhas gerais, pode -se dizer que a combinação de poucos e nenhum decresceu com o desempenho tanto para concluintes como para ingressantes.
Os Gráficos 4.8a e 4.8b apresentam a distribuição das respostas à questão “ Os enunciados das questões, da prova na parte de Formação Geral, estavam claros e objetivos? ”, segundo a Região onde o curso está situado para os estudantes concluintes e ingressantes que participaram do ENADE/2009 na área de ARQUIVOLOGIA.
Em termos regionais, as proporções de ingressantes e de concluintes que consideraram claros e objetivos os enunciados de todas ou da maioria das questões não variaram muito ( no máximo 10 pontos percentuais ).
Entre os concluintes, os valores foram de 71,9 \% no Centro-Oeste a 81,2 \% no Sul.
Já entre os ingressantes, estas proporções variaram de 77,3 \% no Sudeste a 87,5 \% no Norte.
Com relação às respostas “ poucos ” e “ não, nenhum ”, tanto para concluintes quanto para ingressantes, as maiores proporções ocorreram na Região Sudeste, respectivamente 14,3 \% e 10,3 \%. 4.3.2 Componente Específico Os Gráficos 4.9a, 4.9b e 4.9c apresentam a distribuição das respostas à questão “ Os enunciados das questões da prova na parte de formação específica estavam claros e objetivos? ”, segundo quartos de nota final e grupo de estudantes concluintes e ingressantes que participaram do ENADE/2009 na área de ARQUIVOLOGIA.
Com relação aos enunciados das questões da parte do Componente Específico, as opiniões foram também positivas : a grande maioria dos concluintes ( 74,5 \% ) e ingressantes ( 76,6 \% ) considerou claros e objetivos os enunciados de todas ou da maioria das questões.
A proporção de concluintes declarando que todos os enunciados estavam claros e objetivos foi decrescente com o desempenho, indo de 30,0 \% para o primeiro quarto para 16,3 \% para os concluintes com melhor desempenho.
Já a combinação de todas ou da maioria das questões foi crescente com o desempenho, indo de 53,3 \% para os concluintes com pior desempenho a 81,5 \% para o último quarto.
Entre os ingressantes, esta proporção cresceu com o desempenho depois de uma estabilidade nos dois primeiros quartos : em torno de 71,0 \% para os dois piores quartos, alcançando 83,1 \% para o quarto de melhor desempenho.
O número de estudantes que declararam que “ não, nenhum ” dos enunciados das questões da parte do Componente Específico era claro e objetivo, foi maior do que na parte de Formação Geral : 1,8 \% entre os concluintes e 1,2 \% entre os ingressantes.
Os que responderam “ poucos ”, são 5,3 \% entre os concluintes e 7,1 \% entre os ingressantes.
A proporção de estudantes que respondeu que cerca da metade dos enunciados da parte do Componente Específico era claro e objetivo, decresceu com o desempenho, tanto para concluintes quanto para ingressantes.
Os Gráficos 4.10a e 4.10b apresentam a distribuição das respostas à questão “ Os enunciados das questões da prova na parte de formação específica estavam claros e objetivos? ”, segundo a Região onde o curso está situado para os estudantes concluintes e ingressantes que participaram do ENADE/2009 na área de ARQUIVOLOGIA.
Em termos regionais, as proporções de ingressantes e de concluintes que consideraram claros e objetivos os enunciados de todas ou da maioria das questões variaram mais entre os ingressantes ( no máximo 26,1 pontos percentuais ) do que entre os concluintes ( no máximo 13,1 pontos percentuais ).
Entre os concluintes, os valores foram de 70,0 \% no Nordeste a 83,1 \% no Sul.
Com relação às respostas “ poucos ” e “ não, nenhum ” para concluintes a maior proporção ocorreu na Região Sudeste com 9,8 \%. 4.4 Suficiência das informações/instruções fornecidas Os Gráficos 4.11a, 4.11b e 4.11c apresentam a distribuição das respostas à questão “ As informações/instruções fornecidas para a resolução das questões foram suficientes para resolvê -las? ”, segundo quartos de nota final e grupo de estudantes concluintes e ingressantes que participaram do ENADE/2009 na área de ARQUIVOLOGIA.
Quando perguntados sobre as informações/instruções fornecidas para a resolução das questões, quase 83,1 \% dos ingressantes e 84,4 \% dos concluintes de todo o Brasil responderam que estas eram ou excessivas ou suficientes em todas ou na maioria das questões, demonstrando uma avaliação positiva da prova.
Levando -se em conta o desempenho dos participantes, percebeu -se consideráveis diferenças entre as opiniões de estudantes ingressantes e concluintes.
Para os ingressantes, a proporção de estudantes com a opinião supracitada foi crescente com o desempenho.
Entre os concluintes, a proporção vai de 76,6 \% para os de pior desempenho a 87,9 \% para os de melhor desempenho.
Entre os ingressantes, as proporções equivalentes foram, respectivamente, 79,5 \% a 90,2 \%.
O complemento desta opinião, ou seja, os estudantes que responderam somente em algumas ou em nenhuma delas, apresentou uma proporção decrescente com o desempenho para os ingressantes.
A resposta “ não, em nenhuma delas ”, entre os concluintes apareceu somente entre os estudantes de pior performance.
Entre os ingressantes, esta resposta apareceu nos dois extremos : melhores e piores.
Os Gráficos 4.12a e 4.12b apresentam a distribuição das respostas à questão “ As informações/instruções fornecidas para a resolução das questões foram suficientes para resolvê -las? ”, segundo a Região onde o curso está situado para os estudantes concluintes e ingressantes que participaram do ENADE/2009 na área de ARQUIVOLOGIA. 4.5 Dificuldade encontrada ao responder à prova Os Gráficos 4.13a, 4.13b e 4.13c apresentam a distribuição das respostas à questão “ Você se deparou com alguma dificuldade ao responder à prova? ”, segundo quartos de nota final e grupo de estudantes concluintes e ingressantes que participaram do ENADE/2009 na área de ARQUIVOLOGIA.
Perguntados sobre as dificuldades com as quais se depararam ao responder à prova, 40,6 \% dos ingressantes apontaram como principal dificuldade o desconhecimento do conteúdo, já os concluintes apontaram como principal dificuldade a forma diferente de abordagem, 35,0 \%.
O desconhecimento do conteúdo foi apontado somente por 9,3 \% dos concluintes.
A forma diferente de abordagem do conteúdo foi considerada por 25,3 \% dos ingressantes, enquanto que a falta de motivação foi a principal dificuldade para 24,8 \% dos concluintes e para 12,2 \% dos ingressantes.
Entre os ingressantes, os segundo e terceiro quartos apresentaram valores no entorno de 8,0 \%, tendo caído de um valor de 21,7 \% para o quarto de pior desempenho, sendo que, para o grupo de melhor desempenho, este percentual sobe para 14,6 \%. 4.5 Dificuldade encontrada ao responder à prova Os Gráficos 4.13a, 4.13b e 4.13c apresentam a distribuição das respostas à questão “ Você se deparou com alguma dificuldade ao responder à prova? ”, segundo quartos de nota final e grupo de estudantes concluintes e ingressantes que participaram do ENADE/2009 na área de ARQUIVOLOGIA.
Perguntados sobre as dificuldades com as quais se depararam ao responder à prova, 40,6 \% dos ingressantes apontaram como principal dificuldade o desconhecimento do conteúdo, já os concluintes apontaram como principal dificuldade a forma diferente de abordagem, 35,0 \%.
O desconhecimento do conteúdo foi apontado somente por 9,3 \% dos concluintes.
A forma diferente de abordagem do conteúdo foi considerada por 25,3 \% dos ingressantes, enquanto que a falta de motivação foi a principal dificuldade para 24,8 \% dos concluintes e para 12,2 \% dos ingressantes.
Entre os ingressantes, os segundo e terceiro quartos apresentaram valores no entorno de 8,0 \%, tendo caído de um valor de 21,7 \% para o quarto de pior desempenho, sendo que, para o grupo de melhor desempenho, este percentual sobe para 14,6 \%.
Os Gráficos 4.14a e 4.14b apresentam a distribuição das respostas à questão “ Você se deparou com alguma dificuldade ao responder à prova? ”, segundo a Região onde o curso está situado para os estudantes concluintes e ingressantes que participaram do ENADE/2009 na área de ARQUIVOLOGIA.
Em relação às Regiões brasileiras, o percentual de ingressantes que apontaram o desconhecimento do conteúdo como principal dificuldade ao responder à prova variou de 29,6 \% na Região Nordeste, a 80,0 \% na Norte.
Entre os concluintes, tal percentual variou de 3,3 \% na Região Nordeste a 11,5 \%, na Sul.
A forma diferente de abordagem de conteúdo foi citada como maior dificuldade por 21,9 \% dos concluintes na Região Centro-Oeste, chegando a 44,2 \% na Região Sul.
Já entre os ingressantes, os valores foram de 16,0 \% na Região Norte a 28,7 \% na Região Nordeste.
O percentual de concluintes que citaram a falta de motivação como principal dificuldade variou de 15,4 \% na Região Sul a 37,5 \% na Região Centro-Oeste.
Entre os ingressantes, as estatísticas correspondentes foram de 0 \% na Região Norte a 21,3 \%, na Região Centro-Oeste.
Os que declararam que não tiveram qualquer tipo de dificuldade variaram entre os concluintes, de 6,3 \% na Região Centro-Oeste a 21,2 \% na Região Sul.
Entre os ingressantes, este percentual variou de 0,0 \% na Região Norte a 22,2 \% na Região Sul. 4.6 Percepção sobre os conteúdos das questões objetivas da prova Os Gráficos 4.15a, 4.15b e 4.15c apresentam a distribuição das respostas à questão “ Considerando apenas as questões objetivas da prova, você percebeu que :... ”, segundo quartos de nota final e grupo de estudantes concluintes e ingressantes que participaram do ENADE/2009 na área de ARQUIVOLOGIA.
As respostas possíveis para esta questão estavam relacionadas à cobertura do conteúdo das questões objetivas durante o curso.
Como esperado, a proporção de ingressantes ( 52,6 \% ) que declarou que não estudou ainda a maioria desses conteúdos foi bem maior do que entre os concluintes ( 2,6 \% ).
Esta proporção foi decrescente tanto para os concluintes quanto para os ingressantes, como função do desempenho.
Por outro lado, a proporção de estudantes que declarou que estudou e aprendeu todos ou a maioria desses conteúdos, entre os ingressantes foi de 28,9 \% e entre os concluintes 67,9 \%.
Esta proporção foi crescente como função do desempenho, tanto para os concluintes quanto para os ingressantes.
Para os ingressantes, houve uma estabilidade para os dois quartos centrais.
A proporção de estudantes que declarou ter estudado a maioria desses conteúdos, mas não os ter aprendido, foi bem maior entre os concluintes ( 17,2 \% ) do que entre os ingressantes ( 3,7 \% ).
Os Gráficos 4.16a e 4.16b apresentam a distribuição das respostas à questão “ Considerando apenas as questões objetivas da prova, você percebeu que :... ”, para os estudantes concluintes e ingressantes que participaram do ENADE/2009 na área de ARQUIVOLOGIA, segundo a Região onde o curso está situado.
Em relação às Regiões brasileiras, concluintes que afirmaram “ não ter estudado ainda a maioria dos conteúdos ” só ocorreram nas Regiões Nordeste ( 3,3 \% ) e Sudeste ( 4,4 \% ).
Entre os ingressantes, os percentuais variaram de 33,3 \% na Região Centro-Oeste a 88,0 \%, na Norte.
Em todas as Regiões, exceto a Região Norte, a maioria dos concluintes afirmou “ ter estudado e aprendido todos ou muitos dos conteúdos ”, com percentuais que variaram de 66,1 \% na Região Sul a 70,0 \%, na Nordeste.
Para os ingressantes, tal percentual variou de 4,0 \%, na Região Norte, a 44,4 \%, na Centro-Oeste. 4.7 Tempo gasto para concluir a prova Os Gráficos 4.17a, 4.17b e 4.17c apresentam a distribuição das respostas à questão “ Qual foi o tempo gasto para concluir a prova? ”, segundo quartos de nota final e grupo de estudantes concluintes e ingressantes que participaram do ENADE/2009 na área de ARQUIVOLOGIA.
Quando perguntados quanto ao tempo de conclusão da prova, a maioria respondeu entre uma e duas horas, tanto entre os concluintes ( 42,9 \% ) quanto entre os ingressantes ( 39,5 \% ).
A segunda maior incidência de respostas foi entre duas e três horas, igualmente para ambos os grupos : 38,4 \% ( concluintes ) e 37,3 \% ( ingressantes ).
A proporção daqueles que declararam ter usado as quatro horas, mas não ter conseguido terminar foi bem baixa : 0,7 \% entre os ingressantes e 1,3 \% entre os concluintes.
Entre os concluintes esta resposta ficou concentrada quase que exclusivamente entre os estudantes do primeiro quarto, aquele com pior desempenho.
Nem entre os concluintes nem entre os ingressantes percebeu -se uma tendência monotônica do tempo necessário para concluir a prova, como função do desempenho.
Os Gráficos 4.18a e 4.18b apresentam a distribuição das respostas à questão “ Qual foi o tempo gasto para concluir a prova? ”, segundo a Região onde o curso está situado para os estudantes concluintes e ingressantes que participaram do ENADE/2009 na área de ARQUIVOLOGIA.
A proporção de concluintes que gastaram entre uma e duas horas foi maior na Região Sudeste ( 46,4 \% ) e menor na Região Nordeste ( 33,3 \% ).
Para ingressantes, o maior percentual para este tempo ocorreu na Região Norte ( 56,0 \% ) e a menor na Região Nordeste ( 30,2 \% ).
Entre os ingressantes, as Regiões correspondentes foram a Nordeste com tempo mais longo e a Norte com tempo mais curto.

% matched lemmas: laranja
\end{textsample}
